\documentclass[11pt]{article}
\usepackage{amsmath,amssymb,amsthm,hyperref}
\newtheorem{theorem}{Theorem}
\newtheorem{corollary}[theorem]{Corollary}
\title{Linear independence of distinct fair CDFs and a concatenation barrier}
\author{Research proof; independent mathematical reviews completed}
\date{30 September 2026}
\begin{document}
\maketitle

\section{The general CDF theorem}
The following extension of the rank-two calculation was found by the
root collaborator and independently checked by the size-bounds agent.

\begin{theorem}\label{cdf}
Let $n\ge3$ and let independent variables $X_1,\ldots,X_n$ have
continuous distribution functions $F_1,\ldots,F_n$ on a common compact
interval, rescaled to $[0,1]$. Suppose every ordering of every triple
has probability $1/6$. Then:
\begin{enumerate}
\item Every nonzero linear relation among the $F_i$ is supported on
indices having the same distribution function.
\item The distinct functions among $F_1,\ldots,F_n$ are linearly independent.
\item At most one distinct distribution function occurs more than once.
\end{enumerate}
\end{theorem}
\begin{proof}
For distinct $j,k$ set
\[
 d_{kj}=\int_0^1 F_j^2\,dF_k-\frac13.
\]
All Stieltjes integrals below are valid for continuous distribution
functions. Triple fairness gives
$\int F_iF_j\,dF_k=1/3$ when $i,j,k$ are pairwise distinct.
Integration by parts also gives
\[
 \int F_kF_j\,dF_k=\frac12-\frac12\int F_k^2\,dF_j
 =\frac13-\frac12d_{jk}.
\]
Suppose $\sum_i a_iF_i=0$. Evaluating at $1$ gives $\sum_i a_i=0$.
Multiplying by $F_j$ and integrating against $dF_k$, with $j\ne k$,
therefore yields
\[
 0=a_jd_{kj}-\frac12a_kd_{jk}.
\]
Swapping $j,k$ gives $0=a_kd_{jk}-a_jd_{kj}/2$. Consequently
\begin{equation}\label{products}
 a_jd_{kj}=a_kd_{jk}=0\qquad(j\ne k).
\end{equation}

For any two continuous distribution functions $F,G$, integration by
parts yields the nonnegative identity
\begin{equation}\label{pairdivergence}
 \int F^2\,dG+\int G^2\,dF-\frac23
   =\frac12\int(F-G)^2\,d(F+G).
\end{equation}
Its right-hand side is zero exactly when $F=G$.
Indeed write $D=F-G$, $S=F+G$; monotonicity implies
$|D(t)-D(s)|\le S(t)-S(s)$ for $s\le t$.
If, say, $D(t)=d>0$, choose the last $s<t$ with $D(s)=d/2$.
Then $D\ge d/2$ on $[s,t]$ and $S(t)-S(s)\ge d/2$, so the integral
on the right is at least $d^3/16>0$. The case $d<0$ is analogous.

If $a_j,a_k$ are both nonzero, \eqref{products} implies
$d_{kj}=d_{jk}=0$. Equation~\eqref{pairdivergence} then forces $F_j=F_k$.
This proves the first assertion and, by choosing one representative of
each distinct function, the second.
Finally, if two different functions each occurred at least twice, the
relation between two copies of the first function would imply
$d_{kj}=0$ in one direction, and the relation between two copies of the
second would imply $d_{jk}=0$. The same identity would force the two
functions to coincide, a contradiction.
\end{proof}

\begin{corollary}[Exact rank of the block-count matrix]\label{rank}
Let a triple-fair word on $n\ge3$ letters be a concatenation of
permutation-fair blocks with positive multiplicities $v_i^{(g)}$.
Put $C_i=\sum_gv_i^{(g)}$ and let $Q$ be the matrix with columns
\[
 q_i=(v_i^{(g)}/C_i)_{g=1}^s.
\]
If the block-count matrix has rank $r$, then $Q$ has exactly $r$
distinct columns, those columns are linearly independent, and at most
one of them is repeated. In particular, when all normalized letter
profiles are distinct, the block-count vectors span all of
$\mathbb R^n$.
\end{corollary}
\begin{proof}
Replace each block by a time interval of any fixed positive length,
and each die by constant density in that interval with mass
$v_i^{(g)}/C_i$. Independent choices of blocks have exactly the same
probabilities as in the word. Conditional on choosing a common block,
the order of every chosen subset is uniform by the fairness of the
block, matching independent uniform positions in an interval. Thus all
order probabilities, in particular triple fairness, are preserved.
Linear relations between the resulting CDFs are exactly linear relations
between the columns of $Q$: differentiation gives one direction, and
integration with zero initial values gives the other.
Theorem~\ref{cdf} proves the result, since column normalization does not
change the rank of the original count matrix.
\end{proof}

\section{The rank-two calculation and equalization consequences}
The following direct proof is retained because it makes the rank-two
obstruction explicit in two scalar integrals.

\begin{theorem}\label{density}
Let $n\ge3$. Suppose independent random variables $X_1,\ldots,X_n$ on
$[0,1]$ have bounded densities
\[
 f_i(t)=1+a_i h(t),\qquad \int_0^1h(t)\,dt=0,
\]
where $h$ is nonzero on a set of positive measure. If every ordering of
every triple has probability $1/6$, then all but at most one of the
coefficients $a_i$ are equal.
\end{theorem}
\begin{proof}
Put $H(t)=\int_0^t h(x)\,dx$, so the distribution function of $X_i$ is
$F_i(t)=t+a_iH(t)$. Let
\[
 U=\int_0^1 tH(t)\,dt,\qquad V=\int_0^1H(t)^2\,dt>0.
\]
The strict positivity of $V$ follows from absolute continuity of $H$ and
the hypothesis on $h$. For pairwise distinct $i,j,k$, triple fairness
implies
\[
 \Pr(X_i<X_j,\ X_k<X_j)=\int_0^1F_i(t)F_k(t)f_j(t)\,dt=\frac13.
\]
Since $H(0)=H(1)=0$, integration by parts gives
\[
 \int t^2h=-2U,\qquad
 \int tHh=-\frac V2,\qquad \int H^2h=0.
\]
Expanding the preceding probability therefore yields
\begin{equation}\label{identity}
 2U(a_i+a_k-2a_j)
  +V\bigl(2a_i a_k-a_j(a_i+a_k)\bigr)=0.
\end{equation}
Set $b_i=a_i+2U/V$. Equation~\eqref{identity} is exactly
\[
 2b_i b_k=b_j(b_i+b_k).
\]
Subtract the equation with middle index $i$ from the one with middle
index $j$ to obtain
\[
 3b_k(b_i-b_j)=0\qquad(i,j,k\text{ pairwise distinct}).
\]
If all $b_i$ are equal, the conclusion follows. Otherwise choose
$b_i\ne b_j$. Every $b_k$ outside this pair is zero. Since $n\ge3$,
there is such a $k$; the equation with middle index $j$ then gives
$b_i b_j=0$. Hence at most one $b_i$ is nonzero, again proving the claim.
\end{proof}

\begin{theorem}\label{blocks}
Let $W$ be a concatenation of permutation-fair words $W_1,\ldots,W_s$
on the same $n\ge3$ letters, each having positive multiplicities.
Write $v^{(g)}\in\mathbb R_{>0}^n$ for the multiplicity vector of block
$g$, and $C_i=\sum_gv_i^{(g)}$ for the total multiplicities.
Suppose the vectors $v^{(g)}$ span a space of dimension at most two.
If $W$ is triple-fair, then there is an index $e$ such that
\[
 \frac{v_i^{(g)}}{C_i}=\frac{v_j^{(g)}}{C_j}
 \quad\text{for every }g\text{ and all }i,j\ne e.
\]
Thus all but at most one letter must have the same distribution over the
blocks.
\end{theorem}
\begin{proof}
The normalized block-mass vectors
$q^{(g)}=(v_i^{(g)}/C_i)_{i=1}^n$ span a space of dimension at most two,
contain the vector $\mathbf1=\sum_gq^{(g)}$ in their span, and have
strictly positive entries. Choose any positive cell lengths $\lambda_g$
summing to one and represent each die by constant density
$q_i^{(g)}/\lambda_g$ on cell $g$.
This model has exactly the same ordering probabilities as the original
concatenation: the independent block choices have the same masses, and
within each chosen block every restriction of its fair word has uniform
relative order, just as independent uniforms in a common interval do.

One may choose cell lengths so that the density vectors lie on an affine
line through $\mathbf1$. Indeed choose any linear functional $L$ on the
span with $L(\mathbf1)=1$ and $L(q^{(g)})>0$ for all $g$, for example
$L(x)=n^{-1}\sum_i x_i$. Put $\lambda_g=L(q^{(g)})$.
Then each density vector $q^{(g)}/\lambda_g$ belongs to the affine slice
$L(x)=1$. If the span has dimension one, every density is $\mathbf1$ and
the claim is immediate. Otherwise the slice is a line, so
\[
 f_i(t)=1+a_i h(t)
\]
for a fixed vector $a$ and a bounded piecewise-constant scalar $h$.
The vector $a$ is nonconstant because $L(a)=0$ and $a\ne0$; the density
normalizations imply $\int h=0$. If $h$ vanishes identically the claim
is again immediate. Otherwise Theorem~\ref{density} says all but one
$a_i$ are equal, hence all but one density profile and block-mass profile
are identical.
\end{proof}

\begin{corollary}[Two block types]\label{two}
Let $c,d\in\mathbb R_{>0}^n$ and let a word be formed from positive
scaled permutation-fair blocks of count vector $c$ or $d$, with at least
one block of each type. If the concatenation is triple-fair, then the
ratios $c_i/d_i$ are equal for all but at most one index $i$.
\end{corollary}
\begin{proof}
Apply Theorem~\ref{blocks}. On all nonexceptional indices $i,j$, both
$c_i/C_i=c_j/C_j$ and $d_i/C_i=d_j/C_j$ hold, giving the assertion.
\end{proof}

\begin{corollary}[A transposition cannot equalize]
Let $W$ be permutation-fair on $n\ge3$ letters, and suppose letters
$i,j$ have different multiplicities. No concatenation using positive
scaled copies of $W$ and of its relabelling by the transposition $(ij)$,
with at least one of each type, is triple-fair.
\end{corollary}
\begin{proof}
The ratios of the two count vectors are $q,q^{-1},1$ on $i,j$ and any
other letter, where $q>0$ and $q\ne1$. These three ratios are distinct,
contradicting Corollary~\ref{two}.
\end{proof}

\paragraph{Scope.}
The CDF theorem gives a structural rank obstruction for all such
concatenations; its two-type corollary excludes simple transposition
averaging. It gives no quantitative size lower bound for general
equalization and no lower bound on the true equal-composition optimum.
The same argument applies to triple-fair blocks; full permutation
fairness of each block is more than is needed.
\end{document}
